\clearpage
\section{Six unit tests in code}\label{sec:unitcode}
Each listing contains the complete test method. The fixture code below gives
those methods a Student, a Teacher, and a pending application. The tests call
\code{EnrollmentService} directly.

\subsection{Set up the service and its inputs}
\par\smallskip
{\small\href{\repo/blob/5d22aed7f6d4ef1e4c0b0c0dabf2962ba4b47608/src/test/java/com/example/fmschool/EnrollmentServiceUnitTests.java#L33-L37}{\code{EnrollmentServiceUnitTests.java}, lines 33--37}}
\begin{lstlisting}[style=source,language=FMJava,firstnumber=33]
private final EnrollmentRepository enrollments = mock(EnrollmentRepository.class);
private final AccountRepository accounts = mock(AccountRepository.class);
private final TeacherRepository teachers = mock(TeacherRepository.class);
private final StudentRepository students = mock(StudentRepository.class);
private final EnrollmentService service = new EnrollmentService(enrollments, accounts, teachers, students);
\end{lstlisting}

\par\smallskip
{\small\href{\repo/blob/5d22aed7f6d4ef1e4c0b0c0dabf2962ba4b47608/src/test/java/com/example/fmschool/EnrollmentServiceUnitTests.java#L45-L63}{\code{EnrollmentServiceUnitTests.java}, lines 45--63}}
\begin{lstlisting}[style=source,language=FMJava,firstnumber=45]
@BeforeEach
void setUp() {
    teacherAccount = account(2L, "teacher1", Role.TEACHER);
    studentAccount = account(3L, "student1", Role.STUDENT);
    teacher = new Teacher();
    teacher.setId(10L);
    teacher.setAccount(teacherAccount);
    student = student(20L, "Student One", studentAccount);
    application = enrollment(7L, teacher, EnrollmentStatus.PENDING);

    when(enrollments.findStudentId(7L)).thenReturn(Optional.of(20L));
    when(students.findLockedById(20L)).thenReturn(Optional.of(student));
    when(enrollments.findLockedById(7L)).thenReturn(Optional.of(application));
}

@AfterEach
void clearAuthentication() {
    SecurityContextHolder.clearContext();
}
\end{lstlisting}

\par\smallskip
{\small\href{\repo/blob/5d22aed7f6d4ef1e4c0b0c0dabf2962ba4b47608/src/test/java/com/example/fmschool/EnrollmentServiceUnitTests.java#L159-L169}{\code{EnrollmentServiceUnitTests.java}, lines 159--169}}
\begin{lstlisting}[style=source,language=FMJava,firstnumber=159]
private void signInAs(Account account) {
    when(accounts.findByUsername(account.getUsername())).thenReturn(Optional.of(account));
    SecurityContextHolder.getContext().setAuthentication(new TestingAuthenticationToken(
            account.getUsername(), null, "ROLE_" + account.getRole().name()));
}

private static ResponseStatusException assertStatus(HttpStatus status, Runnable action) {
    var error = assertThrows(ResponseStatusException.class, action::run);
    assertEquals(status, error.getStatusCode());
    return error;
}
\end{lstlisting}


\begin{itemize}
\item The four repository mocks stand in for database calls. The service itself
is the real implementation.
\item \code{setUp()} runs before every test. Teacher account 2 owns Teacher
profile 10. Student account 3 owns Student profile 20. Application 7 connects
them and starts as a request from the Student.
\item The three repository stubs let the service follow application 7 to
Student 20 and load the same application object. Mockito supplies those
values when the service calls the repository.
\item \code{signInAs} sets the test identity and stubs account lookup.
The cleanup method clears the security context after each test.
\end{itemize}

\clearpage
\subsection{U1. An acceptance retry must leave the original result intact}
\par\smallskip
{\small\href{\repo/blob/5d22aed7f6d4ef1e4c0b0c0dabf2962ba4b47608/src/test/java/com/example/fmschool/EnrollmentServiceUnitTests.java#L65-L77}{\code{EnrollmentServiceUnitTests.java}, lines 65--77}}
\begin{lstlisting}[style=source,language=FMJava,firstnumber=65]
@Test
void retryingAcceptanceReturnsConflictWithoutSavingAgain() {
    signInAs(teacherAccount);

    service.respond(7L, true);
    assertEquals(EnrollmentStatus.ACCEPTED, application.getStatus());
    verify(enrollments).save(application);
    clearInvocations(enrollments);

    assertStatus(HttpStatus.CONFLICT, () -> service.respond(7L, true));
    assertEquals(EnrollmentStatus.ACCEPTED, application.getStatus());
    verify(enrollments, never()).save(any());
}
\end{lstlisting}


\begin{enumerate}
\item \textbf{Arrange.} Sign in as the Teacher who received the pending
application. Use the shared application with ID 7.
\item \textbf{Act the first time.} \code{respond(7L, true)} accepts the
application. Check that it is accepted and that the service called
\code{save(application)}.
\item \textbf{Separate the two calls.} \code{clearInvocations} clears Mockito's
recorded calls while keeping its stubs and the accepted application object.
The next verification now measures only the retry.
\item \textbf{Act again and assert.} Expect 409. Check the same record still
says accepted and verify that no repository save followed the retry.
\end{enumerate}

\textbf{Why it matters.}
A user can click twice or retry a request after a slow response. The service
must recognize that the first acceptance already took effect. Checking the
exception alone would miss a second write.

\textbf{Find the production guard.}
In \code{EnrollmentService.respond}, the guard requires \code{PENDING}.
The first call changes the status to \code{ACCEPTED}. The retry stops at the
guard before reaching \code{enrollments.save}.

\clearpage
\subsection{U2. A rejected application can reopen as an invitation}
\par\smallskip
{\small\href{\repo/blob/5d22aed7f6d4ef1e4c0b0c0dabf2962ba4b47608/src/test/java/com/example/fmschool/EnrollmentServiceUnitTests.java#L79-L100}{\code{EnrollmentServiceUnitTests.java}, lines 79--100}}
\begin{lstlisting}[style=source,language=FMJava,firstnumber=79]
@Test
void rejectedApplicationReopensAsInvitationWithStudentAsRecipient() {
    signInAs(teacherAccount);
    service.respond(7L, false);
    assertEquals(EnrollmentStatus.REJECTED, application.getStatus());

    when(teachers.findByAccountId(2L)).thenReturn(Optional.of(teacher));
    when(enrollments.findLockedPair(10L, 20L)).thenReturn(Optional.of(application));
    clearInvocations(enrollments);

    service.invite(20L);

    assertEquals(EnrollmentStatus.PENDING, application.getStatus());
    assertEquals(Role.TEACHER, application.getRequestedBy());
    verify(enrollments).save(same(application));
    assertStatus(HttpStatus.FORBIDDEN, () -> service.respond(7L, true));
    assertEquals(EnrollmentStatus.PENDING, application.getStatus());

    signInAs(studentAccount);
    service.respond(7L, true);
    assertEquals(EnrollmentStatus.ACCEPTED, application.getStatus());
}
\end{lstlisting}


\begin{tabularx}{\linewidth}{@{}p{.23\linewidth}p{.22\linewidth}Y@{}}
\toprule
Action & Stored status & Sender and next actor \\
\midrule
Student applies & Pending & Student sends. Teacher may respond. \\
Teacher rejects & Rejected & A new request can reopen the same pair. \\
Teacher invites & Pending & Teacher sends. Student may respond. \\
Student accepts & Accepted & The association is now active. \\
\bottomrule
\end{tabularx}

\begin{itemize}
\item The Teacher first acts as recipient. After rejecting, the mock returns
the same row for that Teacher--Student pair.
\item \code{verify(...).save(same(application))} checks object identity. This
catches replacing the rejected row with a different object.
\item The invitation changes both the status and \code{requestedBy}. The
Teacher's attempt to accept it now gets 403. The Student then accepts it.
\end{itemize}

\textbf{Viva point.}
The owner of the right to respond changes when the sender changes. Merely
checking for \code{PENDING} would miss that rule.

\clearpage
\subsection{U3. Show only eligible Students in the Teacher's picker}
\par\smallskip
{\small\href{\repo/blob/5d22aed7f6d4ef1e4c0b0c0dabf2962ba4b47608/src/test/java/com/example/fmschool/EnrollmentServiceUnitTests.java#L102-L114}{\code{EnrollmentServiceUnitTests.java}, lines 102--114}}
\begin{lstlisting}[style=source,language=FMJava,firstnumber=102]
@Test
void teacherPickerExcludesOccupiedStudentsAndProfilesWithoutAccounts() {
    signInAs(teacherAccount);
    when(teachers.findByAccountId(2L)).thenReturn(Optional.of(teacher));
    Student occupied = student(21L, "Occupied", account(4L, "occupied", Role.STUDENT));
    Student legacy = student(22L, "Legacy", null);
    when(students.findAll()).thenReturn(List.of(student, occupied, legacy));
    when(enrollments.findStudentIdsByStatus(EnrollmentStatus.ACCEPTED)).thenReturn(List.of(21L));

    var board = service.getBoard();

    assertEquals(List.of(new EnrollmentService.StudentOption(20L, "Student One")), board.students());
}
\end{lstlisting}


The repository returns three Students. The test marks profile 21 as already
enrolled and gives profile 22 no account. Comparing the entire returned list
proves both exclusions and confirms that profile 20 remains selectable.

\subsection{U4. Keep pending rows visible when a Student has a course}
\par\smallskip
{\small\href{\repo/blob/5d22aed7f6d4ef1e4c0b0c0dabf2962ba4b47608/src/test/java/com/example/fmschool/EnrollmentServiceUnitTests.java#L116-L131}{\code{EnrollmentServiceUnitTests.java}, lines 116--131}}
\begin{lstlisting}[style=source,language=FMJava,firstnumber=116]
@Test
void enrolledStudentKeepsPendingRowsButCannotPickANewTeacher() {
    signInAs(studentAccount);
    when(students.findByAccountId(3L)).thenReturn(Optional.of(student));
    Teacher otherTeacher = new Teacher();
    otherTeacher.setId(11L);
    Enrollment accepted = enrollment(8L, otherTeacher, EnrollmentStatus.ACCEPTED);
    when(enrollments.findForAccount(3L)).thenReturn(List.of(application, accepted));
    when(teachers.findAll()).thenReturn(List.of(teacher, otherTeacher));

    var board = service.getBoard();

    assertTrue(board.teachers().isEmpty());
    assertEquals(List.of(7L, 8L), board.rows().stream().map(EnrollmentService.Row::id).toList());
    assertEquals(EnrollmentStatus.PENDING, board.rows().getFirst().status());
}
\end{lstlisting}


The existing accepted row empties the Teacher picker. The two row IDs prove that
the board still shows both relationships, and the status assertion confirms the
older request stays pending.

\textbf{The distinction.}
The picker answers whom the user can apply to now. The rows show existing
requests and enrollments. Filtering the picker should leave that history visible.

\clearpage
\subsection{U5. Clearing a mark returns the enrollment to ungraded}
\par\smallskip
{\small\href{\repo/blob/5d22aed7f6d4ef1e4c0b0c0dabf2962ba4b47608/src/test/java/com/example/fmschool/EnrollmentServiceUnitTests.java#L133-L146}{\code{EnrollmentServiceUnitTests.java}, lines 133--146}}
\begin{lstlisting}[style=source,language=FMJava,firstnumber=133]
@Test
void clearingAnExistingMarkRestoresTheUngradedState() {
    signInAs(teacherAccount);
    application.setStatus(EnrollmentStatus.ACCEPTED);
    application.setMarks(new BigDecimal("78.25"));
    EnrolledStudentForm form = new EnrolledStudentForm();
    form.setName(student.getName());
    form.setRoll(student.getRoll());

    service.updateEnrolledStudent(7L, form);

    assertNull(application.getMarks());
    verify(enrollments).save(application);
}
\end{lstlisting}


The accepted enrollment starts with 78.25. The form keeps the profile fields
and leaves marks unset. The service sets the mark to null and saves the
enrollment. This catches a change that would store zero or retain 78.25.

\subsection{U6. A Teacher without a profile cannot invite}
\par\smallskip
{\small\href{\repo/blob/5d22aed7f6d4ef1e4c0b0c0dabf2962ba4b47608/src/test/java/com/example/fmschool/EnrollmentServiceUnitTests.java#L148-L157}{\code{EnrollmentServiceUnitTests.java}, lines 148--157}}
\begin{lstlisting}[style=source,language=FMJava,firstnumber=148]
@Test
void teacherWithAnAccountButNoProfileCannotSendAnInvitation() {
    signInAs(teacherAccount);
    when(teachers.findByAccountId(2L)).thenReturn(Optional.empty());

    var conflict = assertStatus(HttpStatus.CONFLICT, () -> service.invite(20L));

    assertEquals("Complete My Profile before inviting", conflict.getReason());
    verify(enrollments, never()).save(any());
}
\end{lstlisting}


The account exists and has the Teacher role. \code{findByAccountId(2L)} returns
empty, so the service rejects the invite with the profile-completion message.
\code{never().save(any())} verifies there was no enrollment write.

\textbf{Why check the reason?}
Several enrollment rules can return 409. The reason proves that the missing
profile caused the conflict.

\clearpage
\subsection{Read the fixture helpers}
These builders create plain entity objects. Fixed IDs make the scenarios
readable. The integration tests use database-generated IDs instead.

\par\smallskip
{\small\href{\repo/blob/5d22aed7f6d4ef1e4c0b0c0dabf2962ba4b47608/src/test/java/com/example/fmschool/EnrollmentServiceUnitTests.java#L171-L196}{\code{EnrollmentServiceUnitTests.java}, lines 171--196}}
\begin{lstlisting}[style=source,language=FMJava,firstnumber=171]
private Enrollment enrollment(long id, Teacher assignedTeacher, EnrollmentStatus status) {
    Enrollment enrollment = new Enrollment();
    enrollment.setId(id);
    enrollment.setTeacher(assignedTeacher);
    enrollment.setStudent(student);
    enrollment.setRequestedBy(Role.STUDENT);
    enrollment.setStatus(status);
    return enrollment;
}

private static Account account(long id, String username, Role role) {
    Account account = new Account();
    account.setId(id);
    account.setUsername(username);
    account.setRole(role);
    return account;
}

private static Student student(long id, String name, Account account) {
    Student student = new Student();
    ReflectionTestUtils.setField(student, "id", id);
    student.setName(name);
    student.setRoll("00123");
    student.setAccount(account);
    return student;
}
\end{lstlisting}


\begin{itemize}
\item \code{enrollment} connects the supplied Teacher to the shared Student.
Every fixture starts with the Student as sender unless the test changes it.
\item \code{account} supplies the identity the service will resolve during the
test. The Teacher profile's ID and the Teacher account's ID are separate values.
\item \code{student} sets the generated ID through Spring's test utility because
the production entity has no ID setter. Passing null for the account creates
the legacy profile used by U3.
\end{itemize}

To run just these six tests from Windows PowerShell with JDK 25 configured:
\begin{lstlisting}[style=commands]
.\mvnw.cmd '-Dtest=EnrollmentServiceUnitTests' test
\end{lstlisting}

\clearpage
\section{Three integration tests in code}\label{sec:itcode}

\subsection{Start a real application context and use a test database}
\par\smallskip
{\small\href{\repo/blob/5d22aed7f6d4ef1e4c0b0c0dabf2962ba4b47608/src/test/java/com/example/fmschool/EnrollmentPersistenceIT.java#L24-L35}{\code{EnrollmentPersistenceIT.java}, lines 24--35}}
\begin{lstlisting}[style=source,language=FMJava,firstnumber=24]
@SpringBootTest
@ActiveProfiles("integration")
class EnrollmentPersistenceIT {
    private static final String TEST_PASSWORD = "TestPass123!";
    @Autowired private WebApplicationContext context;
    @Autowired private JdbcTemplate db;
    private MockMvc mvc;

    @BeforeEach
    void setUp() {
        mvc = webAppContextSetup(context).apply(springSecurity()).build();
    }
\end{lstlisting}

\par\smallskip
{\small\href{\repo/blob/5d22aed7f6d4ef1e4c0b0c0dabf2962ba4b47608/src/test/resources/application-integration.properties#L1-L5}{\code{application-integration.properties}, lines 1--5}}
\begin{lstlisting}[style=source,language=YAML,firstnumber=1]
# Use a dedicated database for integration tests. Liquibase builds the schema.
spring.docker.compose.enabled=false
spring.datasource.url=jdbc:postgresql://127.0.0.1:55432/fm_school_test
spring.datasource.username=fm_school_test
spring.datasource.password=fm_school_test
\end{lstlisting}


\begin{enumerate}
\item \code{@SpringBootTest} loads the controllers, security, services,
repositories, and Liquibase.
\item \code{@ActiveProfiles("integration")} selects a dedicated PostgreSQL
database at port 55432. The profile disables the app's usual Compose startup.
The test database is supplied by the local disposable container or GitHub Actions.
\item \code{webAppContextSetup} builds MockMvc from the real application context.
\code{springSecurity()} attaches the filter chain. The request runs inside the
test process.
\item \code{JdbcTemplate} checks stored rows after requests finish. These tests
let service transactions commit or roll back normally so the query can check
what actually persisted.
\end{enumerate}

\textbf{Identity and CSRF.}
\code{with(user(...).roles(...))} supplies the authenticated test identity.
The service still resolves that username against a real account. \code{with(csrf())}
supplies the token needed for a write. I3 also sends a login request with the
stored Student's original password.

\textbf{Where the files run.}
Maven Failsafe finds classes ending in \code{IT} during \code{verify}. The
existing startup check and the three scenarios run in this phase.

\clearpage
\subsection{Create fixtures through the API}
\par\smallskip
{\small\href{\repo/blob/5d22aed7f6d4ef1e4c0b0c0dabf2962ba4b47608/src/test/java/com/example/fmschool/EnrollmentPersistenceIT.java#L98-L126}{\code{EnrollmentPersistenceIT.java}, lines 98--126}}
\begin{lstlisting}[style=source,language=FMJava,firstnumber=98]
private Person createPerson(String collection, String name) throws Exception {
    String username = testUsername();
    String body = """
            {"username":"%s","password":"%s","profile":{"name":"%s","roll":"00123","course":"SEPM"}}
            """.formatted(username, TEST_PASSWORD, name);
    mvc.perform(post("/api/" + collection).with(user("admin").roles("ADMIN")).with(csrf())
                    .contentType("application/json").content(body))
            .andExpect(status().isCreated());
    String query = collection.equals("students")
            ? "select s.id from students s join accounts a on s.account_id = a.id where a.username = ?"
            : "select t.id from teachers t join accounts a on t.account_id = a.id where a.username = ?";
    return new Person(db.queryForObject(query, Long.class, username), username);
}

private long apply(Person student, Person teacher) throws Exception {
    mvc.perform(post("/api/enrollments/applications")
                    .with(user(student.username()).roles("STUDENT")).with(csrf())
                    .contentType("application/json")
                    .content("{\"teacherId\":" + teacher.profileId() + "}"))
            .andExpect(status().isCreated());
    return db.queryForObject("select id from enrollments where student_id = ? and teacher_id = ?",
            Long.class, student.profileId(), teacher.profileId());
}

private void accept(Person teacher, long id, int expectedStatus) throws Exception {
    mvc.perform(post("/api/enrollments/{id}/accept", id)
                    .with(user(teacher.username()).roles("TEACHER")).with(csrf()))
            .andExpect(status().is(expectedStatus));
}
\end{lstlisting}


\begin{itemize}
\item \code{createPerson} posts to the Admin API and expects 201. The test
looks up the new profile ID by joining it to the account's unique username.
\item \code{apply} posts the Teacher profile ID as the Student. The following
query retrieves the new enrollment ID for that pair.
\item \code{accept} calls the enrollment route as the specified Teacher and
checks the expected response. The same helper handles successful and
conflicting attempts.
\end{itemize}

The usernames include UUIDs. This keeps fixture names distinct when tests run
again against the same disposable database.

\clearpage
\subsection{I1. A Student can transfer only after leaving the first course}
\par\smallskip
{\small\href{\repo/blob/5d22aed7f6d4ef1e4c0b0c0dabf2962ba4b47608/src/test/java/com/example/fmschool/EnrollmentPersistenceIT.java#L37-L60}{\code{EnrollmentPersistenceIT.java}, lines 37--60}}
\begin{lstlisting}[style=source,language=FMJava,firstnumber=37]
@Test
void studentCanTransferOnlyAfterLeavingTheAcceptedCourse() throws Exception {
    var student = createPerson("students", "Student");
    var first = createPerson("teachers", "First Teacher");
    var second = createPerson("teachers", "Second Teacher");
    long firstRequest = apply(student, first);
    long secondRequest = apply(student, second);

    accept(first, firstRequest, 200);
    accept(second, secondRequest, 409);

    assertEquals("ACCEPTED", enrollmentStatus(firstRequest));
    assertEquals("PENDING", enrollmentStatus(secondRequest));

    mvc.perform(delete("/api/enrollments/{id}", firstRequest)
                    .with(user(student.username()).roles("STUDENT")).with(csrf()))
            .andExpect(status().isNoContent());
    accept(second, secondRequest, 200);

    assertEquals(0, count("select count(*) from enrollments where id = ?", firstRequest));
    assertEquals("ACCEPTED", enrollmentStatus(secondRequest));
    assertEquals(1, count("select count(*) from enrollments where student_id = ? and status = 'ACCEPTED'",
            student.profileId()));
}
\end{lstlisting}


\begin{enumerate}
\item Create one Student and two Teachers. Submit two applications before
accepting either one. This builds two real pending rows.
\item Accept the first application and expect 200. Accepting the second now
returns 409. Queries confirm the first remains accepted and the second remains
pending.
\item Remove the first enrollment as the Student. Expect 204. The second
acceptance now succeeds because the first association is gone.
\item Check that the old row was deleted, the second row is accepted, and the
Student has exactly one accepted enrollment.
\end{enumerate}

\textbf{What this proves.}
The rule is enforced across separate requests and committed database state. A
failed acceptance leaves the pending row available after the Student leaves
the original course.

\clearpage
\subsection{I2. A failed profile save rolls back the account}
\par\smallskip
{\small\href{\repo/blob/5d22aed7f6d4ef1e4c0b0c0dabf2962ba4b47608/src/test/java/com/example/fmschool/EnrollmentPersistenceIT.java#L62-L75}{\code{EnrollmentPersistenceIT.java}, lines 62--75}}
\begin{lstlisting}[style=source,language=FMJava,firstnumber=62]
@Test
void teacherCreationRollsBackTheAccountWhenCourseIsMissing() throws Exception {
    String username = testUsername();
    String body = """
            {"username":"%s","password":"%s","profile":{"name":"Test Teacher"}}
            """.formatted(username, TEST_PASSWORD);

    mvc.perform(post("/api/teachers").with(user("admin").roles("ADMIN")).with(csrf())
                    .contentType("application/json").content(body))
            .andExpect(status().isBadRequest())
            .andExpect(jsonPath("$.message").value("Course is required"));

    assertEquals(0, count("select count(*) from accounts where username = ?", username));
}
\end{lstlisting}


\begin{enumerate}
\item Build valid account data with a unique username and a valid profile name.
Omit only the course. The request can reach the service's role-specific rule.
\item Post the request as Admin with a CSRF token. Check 400 and the message
\code{Course is required}. That confirms the profile save rejected this request.
\item Query the accounts table by username. A count of zero proves the identity
created earlier in the transaction was rolled back.
\end{enumerate}

\textbf{Walk through the call chain.}
\code{TeacherApiController} invokes \code{PeopleApiService.createTeacher}.
The service creates the identity, then calls \code{updateManagedProfile}.
The missing course throws an exception inside the transaction. PostgreSQL
rolls back the account insert, and the exception handler returns 400.

\textbf{Why query the database?}
A controller could return the expected status even if it left an account
behind. The query detects that partial write.

\clearpage
\subsection{I3. Delete a Teacher and keep the Student's login usable}
\par\smallskip
{\small\href{\repo/blob/5d22aed7f6d4ef1e4c0b0c0dabf2962ba4b47608/src/test/java/com/example/fmschool/EnrollmentPersistenceIT.java#L77-L96}{\code{EnrollmentPersistenceIT.java}, lines 77--96}}
\begin{lstlisting}[style=source,language=FMJava,firstnumber=77]
@Test
void deletingATeacherRemovesTheEnrollmentButStudentCanStillSignIn() throws Exception {
    var student = createPerson("students", "Student");
    var teacher = createPerson("teachers", "Teacher");
    long request = apply(student, teacher);
    accept(teacher, request, 200);
    assertEquals("ACCEPTED", enrollmentStatus(request));

    mvc.perform(delete("/api/teachers/{id}", teacher.profileId())
                    .with(user("admin").roles("ADMIN")).with(csrf()))
            .andExpect(status().isNoContent());

    assertEquals(0, count("select count(*) from accounts where username = ?", teacher.username()));
    assertEquals(0, count("select count(*) from teachers where id = ?", teacher.profileId()));
    assertEquals(0, count("select count(*) from enrollments where id = ?", request));
    assertEquals(1, count("select count(*) from students where id = ?", student.profileId()));
    mvc.perform(formLogin("/login").user(student.username()).password(TEST_PASSWORD))
            .andExpect(authenticated().withUsername(student.username()))
            .andExpect(redirectedUrl("/me"));
}
\end{lstlisting}


\begin{enumerate}
\item Create a Teacher and Student, then accept the application. The initial
database check confirms there is a real accepted enrollment to remove.
\item Delete the Teacher profile through the Admin API. Expect 204.
\item Check that the Teacher's account and profile are gone. Check that the
enrollment is gone too. The Student profile must still exist.
\item Submit the Student's original credentials through form login. Check the
authenticated username and the redirect to \code{/me}.
\end{enumerate}

\textbf{Why test login after deletion?}
The profile query proves the Student row survived. The login proves the linked
account, password hash, and Student role still work through Spring Security.
The test therefore covers both storage and a real user action.
